\documentclass[11pt]{article}
\usepackage{fontspec}
\usepackage{amsmath,amssymb}
\usepackage[margin=1in]{geometry}
\usepackage{url}
\begin{document}
\sloppy
\section*{GROK H1624 (2026-10-06 16:24 PT): parent cube 0 at 213 of 315 UNSAT; a re-run bug in the H1612 follow-up is fixed}
\textbf{P versus NP: OPEN.}

\section*{0. What landed since H1612}
Two commits landed after the H1612 note (\texttt{972b0b2e}): \texttt{9f88dac3} adds \path{results/raw/monotone_slice/TWOSATDEC30_q35.txt} and \texttt{75d3587d} adds \path{research/explorations/monotone-slice/2026-10-06_1618_TWOSATDEC30.md}. Both are outside \path{research/agents/grok/} and \path{research/main/}; their content is unread here and nothing from them is cited.

\section*{1. Literature}
None. The claim under test is still this loop's own question: is $L(f,3)<9$ for $f=\texttt{a40ff54e580f5d0a}$?

\section*{2. Plumbing fix (no count changes)}
The H1612 follow-up script fed \texttt{sort -n} output to \texttt{comm}, which needs lexicographic order. As a result it re-ran at 600\,s 23 subcubes that already had 600\,s answers. Every duplicate agreed with the earlier answer (all UNSAT), so no count was wrong, but the compute was wasted. That script was stopped at about 16:24\,PT and replaced by \path{/workspace/scratch_h1624_1006/followup2.sh}, which sorts both lists the same way before \texttt{comm} and runs 3 workers. Its queue was checked against an independent Python recount before launch: both give the same 20 subcubes.

\section*{3. Parent cube 0, second-step split (COMPUTATIONAL, partial)}
Parent cube 0 is $p=(0,0,1,2)$ with 315 subcubes from the H1531 split (PROVED/EXACT as an enumeration of the tight constraints; positive control in H1545). Both scratch copies of \path{split2_r3.py} still have sha256 prefix \texttt{285c61c4}. Recounted by script at 16:25\,PT, taking a 600\,s answer over a 60\,s timeout:
\begin{itemize}
\item \path{probe_c0.out} (60\,s cap, still running) has answered subcubes 0 to 231.
\item \path{rerun600.out} holds 600\,s answers for 29 distinct subcubes, all UNSAT, in 61.7 to 233.1\,s, with no 600\,s timeout. Since H1612 the newly closed ones are 149, 155, 156, 157, 159.
\end{itemize}
Combined: 213 of 315 subcubes UNSAT, 0 SAT. Twenty timed out at 60\,s and are now queued at 600\,s (164, 174, 176, 177, 179, 184, 187, 189, 191, 192, 194, 199, 200, 201, 202, 204, 205, 206, 207, 209), and 82 (232 to 314) are not yet run. Parent cube 0 is not closed.

\section*{4. Standing status}
No new finite witness improves a charged resource in this pass, and no claim is falsified. UNSAT subcubes cannot improve $T$, $S$, or $R$. They would support a lower bound only once every parent cube is closed, and that bound would still rest on the tightness lemma (CANDIDATE). Current bounds for this $f$, carried from earlier notes and not re-derived here: $L(f,2)\le 10$ (COMPUTATIONAL), $L(f,2)=10$ CONDITIONAL on the tightness lemma; $L(f,3)\le 9$ and $L(f,4)\le 8$ (COMPUTATIONAL, H1444). Whether $L(f,3)=8$ stays OPEN. The time-state-representation surface stays CANDIDATE. The SAT-to-charged-model compiler and the product-span obstruction are both OPEN. The D(4,8) semantic-cut bound of 48 stays FALSIFIED/RETIRED.

\section*{Next smallest step}
Let the probe and \path{followup2.sh} finish parent cube 0. If all 315 subcubes end UNSAT, record parent cube 0 as UNSAT (COMPUTATIONAL, solver-certified only, no proof log), then start parent cube 1 at the 600\,s cap directly. Re-simulate any SAT line with independent code before citing it.

\bigskip
\noindent\textbf{Status sentence.} P versus NP remains OPEN: the SAT-to-charged-model compiler and the product-span obstruction are open load-bearing arrows. This pass records that parent cube 0 of the tight $L=8$, $r=3$ split for \texttt{a40ff54e580f5d0a} now has 213 of 315 subcubes UNSAT, 20 timeouts queued at 600\,s, 82 not yet run, and 0 SAT; no charged resource improves, no claim is falsified, and this is not a proof of anything about P versus NP.
\end{document}
